\chapteropen{05}

\begin{cele}
  \item porównać papier, czytnik, telefon i audiobook według tego, co każdy z nich ułatwia, a co utrudnia, bez obietnic co do efektów,
  \item policzyć, ile dłużej trwa ta sama książka w formie audio, i zdecydować, czy ten czas jest czasem wolnym, czy stratą,
  \item rozpoznać, kiedy format wpływa na rozumienie i zapamiętywanie, a kiedy jest tylko kwestią wygody,
  \item ułożyć własną regułę: który format do której książki i jak nie mieszać ich w połowie lektury.
\end{cele}

Format nie jest neutralny. Każdy nośnik zmienia, ile dzieje się w jednej chwili, jak łatwo wrócić do poprzedniego akapitu i jak bardzo męczy się oko albo ucho. Wybór formatu nie zwiększy liczby przeczytanych książek sam z siebie. Może jednak usunąć tarcie, które najczęściej powstrzymuje czytelnika: książka leży w domu, a telefon leży obok.

W tym rozdziale nie będzie rankingu. Będzie rachunek i lista pytań, które pomagają go wykonać. Każdy format ma cechy, które dla jednej osoby są zaletą, a dla innej wadą, i ta różnica zwykle wynika z tego, gdzie czytający spędza czas, a nie z natury samego nośnika.

\section{Papier}

Papierowa książka ma kilka przewag, które łatwo docenić dopiero wtedy, gdy się je straci. Nie wysyła powiadomień. Nie ma baterii, która się kończy w połowie rozdziału. Da się ją zostawić otwartą na stole, a po tygodniu wrócić dokładnie do tej strony. Dobrze znany układ kartki, jej grubość i miejsce, w którym był akapit, pomagają w pamięci przestrzennej: wielu czytelników pamięta, że \glqq to było w dolnej części lewej strony\grqq{}, i z tej wskazówki wraca do treści szybciej niż przez przewijanie.

Papier ma też przewagę, która jest wygodniejsza do obrony niż do zmierzenia: można na nim pisać. Podkreślenie, dopisek na marginesie i zakreślona liczba składają się w zapis myśli, który zostaje z czytelnikiem. Pisanie na marginesie nie zastępuje późniejszej pracy z notatkami (wrócimy do tego w rozdziale o zapamiętywaniu), ale daje materiał do tej pracy.

Słabe strony są równie konkretne. Książka waży, zajmuje miejsce i nie ma funkcji wyszukiwania. Jeśli chcesz odnaleźć definicję sprzed stu stron, przeglądasz spis treści albo kartkujesz. Trzeba ją mieć przy sobie, a w tramwaju, w kolejce czy na przerwie obiadowej zwykle jej nie ma. Kosztuje też więcej za egzemplarz niż wersja cyfrowa, a w wielu domach szybko zajmuje półki, co bywa dla niektórych czytelników presją, a dla innych motywacją.

\section{Czytnik e-booków}

Czytnik z papierem elektronicznym jest najbliższy książce pod względem tego, jak wygląda tekst. Ekran nie świeci własnym światłem, więc czyta się go podobnie jak kartkę, nawet w pełnym słońcu. Bateria wystarcza na tygodnie, a nie na wieczory. Nie ma w nim przeglądarki ani aplikacji do wiadomości, więc sam sprzęt nie zachęca do przełączania się na coś innego.

Dochodzą funkcje, których papier nie ma: regulacja wielkości i kroju czcionki, wyszukiwanie słowa w całej książce, słownik i zakładki, które nie gniotą się w torbie. Dla osób z gorszym wzrokiem albo dla tych, którzy czytają w języku obcym, ta różnica jest realna. Podobnie dla czytelników, którzy chcą przeczytać coś ciężkiego i potrzebują sprawdzać pojęcia bez odkładania lektury.

Wady: czytnik kosztuje, a książki w formacie cyfrowym bywają zabezpieczone tak, że przenosi się je między urządzeniami z trudem. Niektóre pliki PDF z tabelami, wzorami i wieloma kolumnami źle się rozkładają na małym ekranie, co wymaga przełączania widoku, a to psuje ciągłość lektury. Ekran monochromatyczny nie nadaje się do wszystkiego: wykresy w kolorze, zdjęcia i diagramy tracą na czytelności. Zmiana odświeżania strony bywa wolniejsza niż przewinięcie w telefonie, co dla niektórych czytelników jest zaletą, a dla innych przeszkodą.

\section{Telefon i tablet}

Telefon jest najbardziej dostępnym nośnikiem, jaki masz, i to jest jego jedyna pewna przewaga. Książka zawsze jest w kieszeni: w kolejce, w poczekalni, w windzie. Aplikacja zapamiętuje miejsce, w którym skończyłeś, i synchronizuje je między urządzeniami.

Ten sam telefon wysyła powiadomienia, pokazuje wiadomości i ma w zasięgu ręki każdy inny ekran. Krótka przerwa na jeden akapit bardzo często kończy się czterdziestoma minutami przeglądania czegokolwiek. To nie jest kwestia charakteru. Telefon został zaprojektowany tak, by przyciągać uwagę, a książka w aplikacji czytniczej konkuruje z resztą jego funkcji o ten sam ekran.

Jeśli telefon ma być narzędziem do czytania, trzeba ustawić go tak, żeby ta konkurencja była mniejsza: tryb samolotowy albo wyciszenie powiadomień na czas sesji, osobny profil lub osobna przeglądarka do książek, a jeśli to możliwe, tryb skali szarości. Ekran jest mały, więc długie zdania rozbijają się na wiele linijek i łatwiej się wtedy zgubić. Dla dłuższych tekstów telefon sprawdza się przede wszystkim jako uzupełnienie innego formatu, a nie jako jedyny.

\section{Audiobook: tempo i rachunek}

Audiobook zmienia jedną rzecz w sposób, który da się policzyć. Mowa jest wolniejsza niż czytanie, więc ta sama książka trwa dłużej. Przyjmijmy typową książkę popularnonaukową o długości \val{typical.words} słów, co jest zgodne z przyjętą w tej książce miarą typowej pozycji. Dla czytania w tempie 250 słów na minutę:

\begin{przyklad}[Ta sama książka, dwa tempa]
  Czytanie: $80\,000 / 250 = 320$ minut, czyli 5 godzin 20 minut.

  Audio przy 150 słowach na minutę (tempo mowy lektora, \verify{założenie do sprawdzenia dla konkretnego nagrania}): $80\,000 / 150 \approx 533$ minuty, czyli 8 godzin 53 minuty.

  Różnica: około 213 minut, czyli 3 godziny 33 minuty. Stosunek 250 do 150 to $1{,}67$ --- audio zajmuje o dwie trzecie czasu więcej.
\end{przyklad}

Ten rachunek wymaga dwóch uczciwych zastrzeżeń. Pierwsze: tempo czytania 250 słów na minutę to założenie, które zmienia się w zależności od tekstu, celu i zmęczenia, i \verify{dla tekstów popularnonaukowych rzadko jest to wartość stała, a jej zakres zależy od czytelnika}. Drugie: tempo audio zależy od lektora, prędkości odtwarzania i tego, czy słuchasz ze stałym tempem, czy przyspieszasz. Przyspieszenie do 1,5 albo 2 razy skraca czas, ale wpływa na to, ile z mowy da się przetworzyć w jednostce czasu. \verify{wpływ przyspieszenia na rozumienie nagranych treści bywa opisywany jako ograniczony przy wyższych prędkościach, ale wyniki zależą od tekstu i słuchacza}.

Pytanie praktyczne brzmi więc inaczej niż \glqq czy audio jest gorsze\grqq{}. Pytanie brzmi, ile czasu masz, którego w ogóle nie przeznaczasz na czytanie. Spacer, jazda samochodem, sprzątanie, gotowanie, poranna toaleta: to są godziny, w których oczy są zajęte, a uszy nie. Jeśli w tygodniu takich czasów jest na przykład 5 godzin, audiobook może je zamienić na pełną książkę tygodniowo, zamiast tracić je na przeglądanie telefonu. Tabela poniżej, z tego samego generatora co reszta książki, pokazuje, ile typowych książek mieści się w roku przy różnych dziennych minutach i tempach; kolumna 150 słów na minutę odpowiada właśnie słuchaniu.

\generated{books-per-year}

Zasada, która z tego wynika, jest prosta i zarazem kontrowersyjna dla części czytelników:

\begin{zasada}[Książka przesłuchana to przeczytana]
  Jeśli słuchałeś całej książki, a nie tylko jej fragmentu, liczy się ona w rocznym rozliczeniu tak samo jak książka przeczytana oczami. Książka porzucona w połowie się nie liczy, bez względu na to, czy czytałeś ją oczami, czy słuchałeś. Rozliczenie dotyczy kończonych książek, nie formy.
\end{zasada}

Ta zasada ma jedną granicę. Jeśli wiesz, że słuchanie daje ci mniej, na przykład w książkach z wykresami, tabelami, przykładami kodu lub skomplikowanymi argumentami, których nie da się przesłuchać bez stałego powrotu do zdania, to nie jest to nic wstydliwego. To informacja, że dla tej książki lepszy jest inny format.

\begin{falle}[Słuchanie w tle to słuchanie]
  Audiobook odtwarzany podczas zmywania albo prowadzenia rozmowy nie jest czytaniem. Mowa jest zrozumiała, dopóki uwaga ma zapas, i rozmywa się, gdy zajmuje ją coś innego. Uczciwy test: po pięciu minutach zatrzymaj nagranie i powiedz własnymi słowami, co właśnie usłyszałeś. Jeśli nie umiesz, przewiń 30 sekund wstecz. To jest koszt słuchania w tle, który trzeba wliczyć, zanim uzna się, że czas jest \glqq darmowy\grqq{}.
\end{falle}

Czas \glqq darmowy\grqq{} wymaga więc dwóch warunków: nie wymaga od ciebie uwagi, której nie masz, i jest na tyle stały, że da się na nim oprzeć nawyk. Spacer o stałej porze i zmywanie po kolacji spełniają oba. Przejazd tramwajem w godzinach szczytu już częściej nie.

\section{Biblioteka i e-wypożyczalnie}

Biblioteka jest najtańszym formatem w historii czytania, bo książka kosztuje tyle, ile wynosi wypożyczenie. Ma też cechę, której żaden zakup nie daje: termin zwrotu. Termin jest uciążliwy, ale dla części czytelników działa jak zewnętrzne zobowiązanie, a nie wewnętrzna obietnica. Książka, która ma zwrot za trzy tygodnie, dostaje pierwszeństwo przed tą, która stoi na półce i będzie czekać.

Wady są równie oczywiste. Kolejka na popularne tytuły bywa długa i nie zawsze da się przewidzieć, kiedy książka będzie dostępna. Dostępność zależy od biblioteki, a więc \verify{zasady wypożyczania, limity i długość kolejek różnią się między bibliotekami i e-wypożyczalniami, więc przed planowaniem sprawdź je u konkretnego dostawcy}. Książka, która przychodzi w kolejce, zwykle przychodzi wtedy, gdy masz już inny plan, i to jest najczęstszy powód, dla którego wypożyczona książka wraca nieprzeczytana.

Praktyczna rada: nie wypożyczaj od razu pięciu tytułów. Wypożycz dwa, z których jeden już wiesz, że przeczytasz w terminie, i jeden, który chcesz sprawdzić. Przy każdej książce z kolejki zapisz datę, do której realnie ją skończysz, a nie datę, do której musisz oddać. Różnica między tymi dwiema datami jest zwykle kilka tygodni i to ona decyduje, czy wypożyczenie przyniesie efekt.

\section{Światło i zmęczenie}

Zmęczenie oczu przy czytaniu jest realne, ale jego przyczyny są mniej oczywiste, niż zakłada wiele poradników. Jedna sprawa jest pewna z doświadczenia: przy słabym oświetleniu i za małej czcionce czytanie męczy szybciej. Druga sprawa, do której trzeba podchodzić ostrożnie: wiele ogólnych twierdzeń o ekranach, światle niebieskim i liczbie mrugnięć wymaga pomiaru w konkretnych warunkach, a nie wiary. \verify{dane o częstości mrugania przy czytaniu z ekranu i z papieru oraz o wpływie światła ekranu na sen różnią się w zależności od badania, więc nie należy traktować ich jako ustalonego faktu}.

Co da się zrobić bez czekania na wyniki badań? Dobrać jasność tak, żeby nie świecić w oczy w ciemnym pokoju. Ustawić czcionkę tak, żeby nie trzeba było przybliżać ekranu. Czytać w miejscu, w którym światło pada na stronę z przodu lub z boku, a nie zza pleców. Przy dłuższej lekturze robić krótkie przerwy: to jest zasada wyboru, a nie wynik badania, i warto ją traktować jako nawyk, który sprawdzisz na własnym zmęczeniu.

Zmęczenie oczu to też zmęczenie uwagi. Książka czytana po całym dniu przed ekranem zwykle wymaga więcej wysiłku niż ta sama książka czytana rano. Jeśli wieczorem przy każdej próbie czytania zasypiasz po dwóch stronach, problemem nie jest format, tylko pora. Wtedy warto przesunąć czytanie na inny moment dnia albo wybrać format, który nie wymaga oczu w tym samym stopniu, czyli audio.

\section{Łączenie formatów}

Ta sama książka może istnieć w dwóch formach jednocześnie: papierowej i audio. Ma to sens w dwóch przypadkach. Pierwszy: czytasz tekst przy biurku, a w dojazdach słuchasz tej samej książki, żeby nie tracić ciągłości. Drugi: słuchasz całości, a do jednego rozdziału wracasz papierem, bo tam są tabele i wzory, które wymagają zatrzymania się.

Taki podział jest korzystny, jeśli każdy format ma swoje zadanie. Audio zapewnia ciągłość i dostęp w chwilach, w których papier jest niedostępny. Papier zapewnia precyzję tam, gdzie trzeba wrócić do zdania, zatrzymać się na przykładzie albo zapisać coś ręcznie. Łączenie formatów przestaje działać wtedy, gdy oba służą do tego samego: wtedy czytasz to samo dwa razy i nie wiesz, który z tych dwóch kontaktów był naprawdę przeczytaniem.

Kłopot pojawia się, gdy zmieniasz format w połowie książki, bez decyzji. Zakładka w papierze i pozycja w aplikacji audio nie odpowiadają sobie wprost, a przełączenie wymaga szukania miejsca, które ci przerwało. Każde takie przełączenie kosztuje chwilę zagubienia, a te chwile kumulują się w porzuconą książkę. Stąd prosta reguła:

\begin{zasada}[Jeden format na książkę]
  Wybierz format przed pierwszą stroną i trzymaj się go do końca, chyba że zmieniasz na świadomie wybrany drugi, o którym decydujesz z góry, np. audio w dojazdach i papier w weekend. Nie przełączaj się w środku rozdziału, bo to właśnie tam gubi się większość czytelników.
\end{zasada}

\begin{uwaga}[Format a książka]
  Nie każda książka wymaga tego samego formatu. Powieść z wieloma postaciami często czyta się dobrze uchem, a książka z kodem lub tabelami ma swoje miejsce na papierze lub ekranie. Zanim zaczniesz, sprawdź, czego ta konkretna książka wymaga, a nie tylko, co ci wygodniej.
\end{uwaga}

Na koniec: wybór formatu jest decyzją o cenie tarcia, a nie o jakości czytania. Jeśli zauważysz, że książka nie czyta się od tygodnia, zapytaj, który element tarcia jest winien: czy to daleko położona książka, bateria, czy ekran telefonu. Zwykle da się to naprawić jedną zmianą, a nie zmianą całego nawyku.

\begin{cwiczenia}
  \exitem{Ile godzin zajmie audio książki o długości 80\,000 słów przy 150 słowach na minutę? Podaj wynik w godzinach i minutach.}{$80\,000 / 150 \approx 533{,}3$ minuty, czyli 8 godzin 53 minuty (około).}
  \exitem{Ta sama książka przy czytaniu w tempie 250 słów na minutę zajmuje 5 godzin 20 minut. Ile dłużej trwa audio i jaki jest stosunek czasów?}{Różnica to około 3 godziny 33 minuty (213 minut). Stosunek to $533{,}3 / 320 \approx 1{,}67$.}
  \exitem{Masz w tygodniu 5 godzin spacerów i zmywania, w których oczy są zajęte. Po ile książek o długości 80\,000 słów miesięcznie możesz tak przesłuchać, zakładając 150 słów na minutę?}{5 godzin tygodniowo to 300 minut, czyli $300 \cdot 150 = 45\,000$ słów tygodniowo. Przy 80\,000 słów na książkę to $45\,000 / 80\,000 \approx 0{,}56$ książki tygodniowo, czyli około 2{,}25 książki miesięcznie (4 tygodnie), przy założeniu, że cały czas idzie na słuchanie.}
  \exitem{Podaj jedną sytuację, w której słuchanie w tle nie liczy się jako czytanie, i test, który to sprawdza.}{Słuchanie podczas prowadzenia rozmowy albo przy zmywaniu, z dużym zaangażowaniem uwagi w inną czynność. Test: po pięciu minutach zatrzymaj nagranie i streść własnymi słowami, co usłyszałeś.}
  \exitem{Wymień dwie rzeczy, które papier daje lepiej niż czytnik, i jedną, której papier nie daje wcale.}{Lepiej niż czytnik: pamięć przestrzenna położenia tekstu na stronie oraz pisanie na marginesie. Czego papier nie daje: wyszukiwania słowa w całej książce (ani zmiany rozmiaru czcionki).}
  \exitem{Co trzeba sprawdzić przed wypożyczeniem książki z kolejki, zanim zadeklarujesz, że ją przeczytasz?}{Długość kolejki i datę, do której realnie skończysz lekturę, a nie termin zwrotu; zasady wypożyczania u konkretnego dostawcy.}
  \exitem{Dlaczego zmiana formatu w połowie rozdziału jest bardziej kosztowna, niż wygląda?}{Trzeba szukać miejsca, w którym przerwano, i odtworzyć kontekst; każda taka zmiana to chwila zagubienia, a kilka z nich kończy się porzuceniem książki.}
\end{cwiczenia}

\zadanie{Wybierz jedną książkę, którą masz w planie na najbliższe dwa tygodnie. Zapisz, w jakim formacie ją zaczynasz, i zapisz w jednym zdaniu, dlaczego. Przez dwa tygodnie nie zmieniaj formatu. Miara sukcesu: na koniec drugiego tygodnia wiesz, ile stron lub minut (dla audio) przeczytałeś, i zapisujesz liczbę dni, w których czytanie się odbyło. Jeśli zmieniłeś format, zapisz, w którym momencie i dlaczego.}
